% =====================================================================
%  Knowledge-Graph Reconfiguration:
%  Catalytic Path-Introduction in Finite Contact Graphs
%
%  A graph-theoretic account of how a finite agent's reachable
%  structure is altered by minimal committed reshuffles, and the
%  consequences thereof.
%
%  Author: Kundai Farai Sachikonye
%          Technical University of Munich
% =====================================================================
\documentclass[11pt,a4paper]{article}

\usepackage[utf8]{inputenc}
\usepackage[T1]{fontenc}
\usepackage{lmodern}
\usepackage[a4paper,margin=2.5cm]{geometry}
\usepackage{amsmath,amssymb,amsthm,mathtools}
\usepackage{bm}
\usepackage{mathrsfs}
\usepackage{booktabs}
\usepackage{enumitem}
\usepackage{microtype}
\usepackage[round,sort&compress]{natbib}
\usepackage[colorlinks=true,linkcolor=blue!45!black,
            citecolor=blue!45!black,urlcolor=blue!45!black]{hyperref}
\usepackage{cleveref}

% ---- Theorem environments ----
\newtheoremstyle{thmstyle}{\topsep}{\topsep}{\itshape}{}{\bfseries}{.}{.5em}{}
\theoremstyle{thmstyle}
\newtheorem{theorem}{Theorem}[section]
\newtheorem{lemma}[theorem]{Lemma}
\newtheorem{proposition}[theorem]{Proposition}
\newtheorem{corollary}[theorem]{Corollary}
\theoremstyle{definition}
\newtheorem{definition}[theorem]{Definition}
\newtheorem{axiom}{Premise}
\newtheorem{construction}[theorem]{Construction}
\theoremstyle{remark}
\newtheorem{remark}[theorem]{Remark}
\newtheorem{notation}[theorem]{Notation}

\crefname{axiom}{Premise}{Premises}
\Crefname{axiom}{Premise}{Premises}

% ---- Notation ----
\newcommand{\R}{\mathbb{R}}
\newcommand{\N}{\mathbb{N}}
\newcommand{\Gph}{\Gamma}
\newcommand{\V}{V}
\newcommand{\E}{E}
\newcommand{\wt}{w}
\newcommand{\floorc}{\beta}
\newcommand{\Nbr}{N}
\newcommand{\cut}{\partial}
\newcommand{\Res}{\rho}
\newcommand{\Inv}{\iota}            % invariant cell
\newcommand{\Walk}{W}
\newcommand{\rec}{M}                % committed count
\newcommand{\hol}{\eta}            % holonomy
\newcommand{\reshuffle}{\sigma}    % a reshuffle
\newcommand{\abs}[1]{\lvert#1\rvert}

\title{\textbf{Knowledge-Graph Reconfiguration:\\[2pt]
Catalytic Path-Introduction in Finite Contact Graphs}}

\author{%
Kundai Farai Sachikonye\\
\small Technical University of Munich\\
\small \texttt{kundai.sachikonye@wzw.tum.de}}

\date{\today}

% =====================================================================
\begin{document}
\maketitle

% =====================================================================
\begin{abstract}
\noindent
We study how the reachable structure of a finite weighted graph changes
under minimal committed edits, and what such edits can and cannot do. The
objects are elementary: a finite weighted graph whose vertices are the
distinguishable parts an agent has resolved, whose edges carry the cost of
separating one part from another, and whose distinguished cells we call
\emph{invariants}---multiply-reachable subgraphs preserved under
relabelling. We take a single premise: the graph is finite and its edge
weights are bounded below by a positive floor. From it we derive: (T0) a
strictly positive floor on every separating cut; (T1) that an invariant is
fixed by its complement, never by a single vertex, hence is a region and
not a point; (T2) a \emph{holonomy} functional on cycles that vanishes
exactly on consistent reachability and is nonzero on inconsistent
reachability; (T3) that the minimal committed edit---a
\emph{reconfiguration unit}---introduces or reweights exactly one
floor-crossing path and increments a monotone committed count, hence is
irreversible; (T4) that a reconfiguration is a \emph{catalyst}: it alters
which paths reach an invariant without being itself a vertex or a path, and
is therefore not extractable from the reconfigured graph; (T5) that a
finite expression of a reconfiguration always carries removable slack, so
no finite expression is identical to the reconfiguration it specifies; and
(T6) a four-fold classification of reconfigurations by two binary
invariants---whether the introduced path closes holonomy, and whether
alternate paths to the target cell survive---together with the result that
a reconfiguration whose stated endpoints close holonomy may still route
through a cell whose holonomy does not close (a \emph{bridge}), detectable
only by auditing the route and not the endpoints. We close with a Master
Theorem: every consequence T0--T6 is forced for finite floored graphs and
fails when finiteness or the floor is dropped. All claims are proved of the
abstract graph; interpretive readings are confined to remarks on which no
theorem depends.
\end{abstract}

\noindent\textbf{Keywords:} finite weighted graph; edge-weight floor;
reachability; cycle holonomy; graph reconfiguration; catalytic edit;
monotone committed count; invariant cell; path-multiplicity.

\tableofcontents

% =====================================================================
\section{Introduction}
\label{sec:intro}
% =====================================================================

\subsection{The question}

A finite agent resolves its surroundings into finitely many distinguishable
parts and the costs of telling them apart. Any such resolution is a finite
weighted graph. This paper asks one structural question about such graphs:
\emph{what happens to the set of paths reaching a distinguished cell when
the graph is minimally edited, and what can such an edit do?} We answer with
a chain of consequences forced by a single premise---that the graph is
finite and positively weighted.

The development is deliberately confined to graph theory. We define a
\emph{thing}, a \emph{resolution} (``knowledge''), an \emph{invariant}
(``truth''), and a \emph{reconfiguration} (``the propagation of an idea'')
entirely as structures on a finite weighted graph, and prove their
properties as theorems. The everyday words are offered only in remarks,
explicitly marked, on which no theorem depends. This keeps the results
where they belong: claims about finite weighted graphs, true or false by
the standards of graph theory alone.

\subsection{Method: conditional consequence}

The paper proves a conditional. \emph{If} a finite floored contact graph
exists, \emph{then} a determinate chain of structural facts about its
reconfigurations is forced. We assume nothing about what the vertices
\emph{are}; we assume only that they are finite in number and positively
separated, and we read off what their reconfigurations cannot lack. The
deductions are combinatorial and order-theoretic throughout.

\subsection{The single premise}

\begin{axiom}[Finiteness with a floor]
\label{ax:finite}
A \emph{contact graph} is a finite weighted graph $\Gph=(\V,\E,\wt)$ with
$\abs{\V}<\infty$, connected, in which every edge weight is bounded below by
a positive constant: there is $\floorc>0$ with $\wt(e)\ge\floorc$ for all
$e\in\E$.
\end{axiom}

The premise is almost nothing: a finite graph whose edges cost something.
Its consequences are the subject of the paper.

\subsection{The chain of consequences}

Under \Cref{ax:finite} we prove:

\begin{enumerate}[label=\textbf{(T\arabic*)},leftmargin=3em,start=0]
\item \emph{The floor} (\Cref{sec:floor}): every separating cut has weight
$\ge\floorc>0$; no two parts are separated at zero cost.
\item \emph{Invariants are regions} (\Cref{sec:invariant}): a
multiply-reachable cell is fixed by its complement, realised by no single
vertex, and preserved under relabelling.
\item \emph{Holonomy} (\Cref{sec:holonomy}): a functional on cycles that is
zero on consistent reachability and nonzero on inconsistent reachability.
\item \emph{The reconfiguration unit} (\Cref{sec:unit}): the minimal
committed edit introduces one floor-crossing path and increments a monotone
count; it is irreversible.
\item \emph{Reconfiguration is catalytic} (\Cref{sec:catalyst}): an edit
alters reachability without being a vertex or a path, hence is not
extractable from the result.
\item \emph{Expressive slack} (\Cref{sec:slack}): every finite specification
of a reconfiguration carries removable slack, so no specification is
identical to what it specifies.
\item \emph{The four modes and the bridge} (\Cref{sec:modes}):
reconfigurations classify by (closes-holonomy?) $\times$
(alternates-survive?), and an edit with holonomy-closing endpoints may route
through a non-closing cell.
\end{enumerate}

\Cref{sec:master} proves these are forced for finite floored graphs and fail
otherwise.

% =====================================================================
\section{The contact graph and its cells}
\label{sec:setting}
% =====================================================================

Throughout, $\Gph=(\V,\E,\wt)$ is a contact graph (\Cref{ax:finite}): finite,
connected, weighted, with floor $\floorc>0$. Graphs are simple unless
stated.

\begin{definition}[Cut]
\label{def:cut}
For $S,T\subseteq\V$ disjoint, the \emph{cut} is
$\cut(S,T)=\{uv\in\E:u\in S,\,v\in T\}$ with weight
$\wt(\cut(S,T))=\sum_{e\in\cut(S,T)}\wt(e)$. We write $\cut(S):=\cut(S,\V\setminus S)$.
\end{definition}

\begin{definition}[Cell; induced subgraph]
\label{def:cell}
A \emph{cell} is a nonempty $U\subseteq\V$; its \emph{induced subgraph}
$\Gph[U]$ has vertex set $U$ and the edges of $\E$ with both ends in $U$. A
cell is \emph{proper} if $\varnothing\neq U\subsetneq\V$.
\end{definition}

\begin{notation}[Vocabulary, on which nothing depends]
\label{not:vocab}
A reader may wish to read a vertex as a \emph{distinguished part}, an edge
weight as the \emph{cost of telling two parts apart}, and a cell as
\emph{something an agent has resolved}. We use only the graph terms; the
readings are orientation. In particular ``thing,'' ``agent,'' ``knowledge,''
``truth,'' ``idea,'' and ``propagation'' name graph structures defined
below and carry no commitment beyond those definitions.
\end{notation}

% =====================================================================
\section{T0: the floor}
\label{sec:floor}
% =====================================================================

\begin{theorem}[Floor]
\label{thm:floor}
Every proper cell $U\subsetneq\V$ has $\wt(\cut(U))\ge\floorc>0$. No proper
cell is separated from its complement at zero cost.
\end{theorem}

\begin{proof}
Since $\Gph$ is connected and $U$ is a nonempty proper subset, $\cut(U)\neq
\varnothing$ (else $\Gph$ disconnects). Then $\wt(\cut(U))\ge\wt(e^\ast)\ge
\floorc>0$ for any $e^\ast\in\cut(U)$, using $\wt(e)\ge\floorc$.
\end{proof}

\begin{remark}[Reading]
\label{rem:floor-read}
Read: ``no distinction is free.'' Telling any part from the rest costs at
least $\floorc$. Nothing below depends on this reading; T0 is the statement
that a finite positively-weighted graph has no zero-weight cut.
\end{remark}

% =====================================================================
\section{T1: invariants are regions, not points}
\label{sec:invariant}
% =====================================================================

We single out the cells that are stable under relabelling and reached by
more than one route. These are the objects we will call invariants.

\begin{definition}[Reachability and path-multiplicity]
\label{def:paths}
A \emph{path} from $u$ to a cell $U$ is a walk $u=v_0,v_1,\dots,v_k$ with
$v_k\in U$ and each $v_iv_{i+1}\in\E$. The cell $U$ is
\emph{$m$-reachable from $u$} if there are $m$ internally edge-disjoint such
paths. We write $\pi(u,U)$ for the maximum number of edge-disjoint paths
from $u$ to $U$.
\end{definition}

\begin{definition}[Invariant cell]
\label{def:invariant}
A cell $\Inv\subseteq\V$ is an \emph{invariant} if (i) it is
\emph{multiply-reachable}: $\pi(u,\Inv)\ge2$ for every $u\notin\Inv$ in the
same component; and (ii) it is \emph{relabelling-stable}: for every weighted
automorphism $\phi$ of $\Gph$, $\phi(\Inv)$ is again an invariant with
$\wt(\cut(\phi(\Inv)))=\wt(\cut(\Inv))$.
\end{definition}

\begin{theorem}[Invariants are regions]
\label{thm:region}
Let $\Inv$ be an invariant. Then:
\begin{enumerate}[label=\textup{(\roman*)},leftmargin=2.2em]
\item \textbf{Positive boundary.} $\wt(\cut(\Inv))\ge\floorc>0$.
\item \textbf{Not a point.} If $\abs{\V}\ge3$, no single vertex $v$ realises
$\Inv$ in the sense $\Inv=\{v\}$ with $\pi(u,\{v\})\ge2$ forced for all $u$;
generically $\abs{\Inv}\ge2$ or $\Inv$ is reached by edge-disjoint paths
through distinct boundary vertices. The invariant is borne by a region of
its boundary, not by one vertex.
\item \textbf{Complement-fixed.} $\Inv$ is determined by its complement:
$\Inv=\V\setminus(\V\setminus\Inv)$, and no rule naming a single vertex of
$\Inv$ in advance is needed to fix it.
\end{enumerate}
\end{theorem}

\begin{proof}
(i) is \Cref{thm:floor} applied to $U=\Inv$. (ii): multiple reachability
($\pi\ge2$) requires at least two edge-disjoint paths into $\Inv$; by
Menger's theorem \citep{diestel2017} the minimum vertex cut separating $u$
from $\Inv$ has size $\ge2$, so no single boundary vertex carries all routes;
hence the ``locus'' of $\Inv$'s reachability is a set of $\ge2$ boundary
vertices, a region. (iii): complementation $U\mapsto\V\setminus U$ is an
involution on subsets of a finite set, so specifying $\V\setminus\Inv$
fixes $\Inv$ uniquely with no further datum.
\end{proof}

\begin{remark}[Reading]
\label{rem:truth-read}
Read an invariant as a ``truth'': a cell reached by many independent routes,
stable under re-encoding, fixed by what it is not, and never collapsible to
a single point. Multiple-reachability is the graph form of
``independently anchored''; relabelling-stability is the graph form of
``the same regardless of how it is described.'' No theorem below depends on
this reading.
\end{remark}

% =====================================================================
\section{T2: holonomy on cycles}
\label{sec:holonomy}
% =====================================================================

We now equip the graph with a consistency functional on cycles. It will
distinguish reachability that closes up from reachability that does not.

\begin{definition}[Transport and holonomy]
\label{def:holonomy}
Let $g:\E\to A$ assign to each edge an element of an abelian group $(A,+)$,
the \emph{transport}, with $g(uv)=-g(vu)$. For a cycle
$c=(v_0,v_1,\dots,v_k=v_0)$, the \emph{holonomy} of $c$ is
\[
  \hol(c)\;=\;\sum_{i=0}^{k-1} g(v_iv_{i+1})\;\in\;A.
\]
The transport is \emph{consistent} if $\hol(c)=0$ for every cycle $c$.
\end{definition}

\begin{proposition}[Consistency is path-independence]
\label{prop:pathindep}
A transport $g$ is consistent if and only if, for all $u,v$ and all paths
$P,P'$ from $u$ to $v$, $\sum_{e\in P}g(e)=\sum_{e\in P'}g(e)$.
\end{proposition}

\begin{proof}
If $g$ is consistent, then $P$ followed by the reverse of $P'$ is a cycle of
holonomy $0$, giving $\sum_P g=\sum_{P'} g$. Conversely, any cycle is a path
from $v_0$ to itself, and path-independence forces its sum to equal the
empty-path sum, $0$.
\end{proof}

\begin{theorem}[Holonomy detects inconsistent reachability]
\label{thm:holonomy}
Let $\Inv$ be an invariant and $g$ a transport. The reachability of $\Inv$
is \emph{consistent} (all routes to $\Inv$ agree under $g$) if and only if
every cycle meeting $\Inv$ has $\hol=0$. If some cycle through the routes to
$\Inv$ has $\hol\neq0$, the routes disagree: there exist two paths to $\Inv$
whose transports differ.
\end{theorem}

\begin{proof}
Immediate from \Cref{prop:pathindep}: agreement of all routes to $\Inv$ is
path-independence on the paths into $\Inv$, equivalent to vanishing holonomy
on cycles formed by pairs of such paths. A nonzero cycle holonomy exhibits,
by construction, two paths with unequal transport sums.
\end{proof}

\begin{remark}[Reading]
\label{rem:holonomy-read}
Read $\hol(c)=0$ as ``the accounts you reach by going around different ways
agree.'' A true invariant closes all its cycles (zero holonomy): every
independent route to it tells the same story. A claim that passes a check
on one route but fails when you go around another way is exactly a cycle of
nonzero holonomy. Nothing below depends on the reading; T2 is a statement
about group-valued transports on cycles.
\end{remark}

% =====================================================================
\section{T3: the reconfiguration unit}
\label{sec:unit}
% =====================================================================

We now define the minimal committed edit of a contact graph and prove it is
quantised and irreversible.

\begin{definition}[Reconfiguration]
\label{def:reconfig}
A \emph{reconfiguration} of $\Gph=(\V,\E,\wt)$ is an edit producing
$\Gph'=(\V,\E',\wt')$ on the same vertex set by adding edges, deleting edges,
or changing weights, such that $\Gph'$ is again a contact graph (finite,
connected, floor $\ge\floorc$). A reconfiguration is \emph{path-introducing}
for a cell $U$ if $\pi'(u,U)>\pi(u,U)$ for some $u$ (a new edge-disjoint
route to $U$ appears).
\end{definition}

\begin{definition}[Committed count]
\label{def:count}
Fix an ordering of committed edits. The \emph{committed count} $\rec$ after a
finite sequence of reconfigurations is the number of edits performed (each
add, delete, or reweight that changes $\Gph$ counts one). A
reconfiguration is \emph{committed} when it is recorded; the record only
grows.
\end{definition}

\begin{definition}[Reconfiguration unit]
\label{def:unit}
A \emph{reconfiguration unit} \textup{(}for a cell $U$\textup{)} is a
path-introducing reconfiguration that is \emph{minimal}: it introduces
exactly one new floor-crossing path to $U$ (raising $\pi(\cdot,U)$ by one
along a single route), with no proper sub-edit also path-introducing for $U$.
\end{definition}

\begin{theorem}[The unit is floored and quantised]
\label{thm:unit-floor}
Every reconfiguration unit for $U$ deposits boundary content $\ge\floorc>0$
along its introduced route, and raises $\pi(\cdot,U)$ by exactly one. There
is no path-introducing reconfiguration of content below $\floorc$.
\end{theorem}

\begin{proof}
A new floor-crossing route to $U$ contains at least one edge crossing
$\cut(U')$ for some $U'\supseteq U$; by \Cref{thm:floor} that edge has weight
$\ge\floorc$, so the introduced route carries content $\ge\floorc$. An edit
that introduces a route entirely of edges below $\floorc$ would violate the
floor of $\Gph'$ as a contact graph (\Cref{def:reconfig}); hence no such
sub-floor path-introduction exists. Minimality (\Cref{def:unit}) gives
exactly one new edge-disjoint route, so $\pi(\cdot,U)$ rises by one.
\end{proof}

\begin{theorem}[Non-return]
\label{thm:nonreturn}
The committed count $\rec$ is strictly monotone: each committed
reconfiguration increases $\rec$ by one, and no committed reconfiguration is
undone by a further reconfiguration without itself being a further committed
edit (which increases $\rec$ again). Consequently no finite sequence returns
the system to a prior committed state: a configuration revisited as a
\emph{graph} is never revisited as a \emph{state} $(\Gph,\rec)$.
\end{theorem}

\begin{proof}
Each committed edit appends one to the record (\Cref{def:count}), so $\rec$
strictly increases. ``Undoing'' edit $k$ is itself an edit $k'$ with
$\rec_{k'}=\rec_k+1>\rec_k$; the resulting graph may equal an earlier graph,
but the state $(\Gph,\rec)$ differs in its second component. Hence for
$i<j$, the states at steps $i$ and $j$ are distinct even when the graphs
coincide.
\end{proof}

\begin{remark}[Reading]
\label{rem:unit-read}
Read a reconfiguration unit as ``an idea in unit form'': the minimal
committed change to an agent's resolved structure that actually introduces a
new route to an invariant. \Cref{thm:unit-floor} reads ``an idea must carry
at least a floor's worth of new uncertainty to move the graph at all''; below
the floor nothing is introduced. \Cref{thm:nonreturn} reads ``a received
reconfiguration cannot be un-received''---the graph can only be reconfigured
further, never restored to the pre-edit state, because the committed count
is part of the state. No theorem depends on these readings.
\end{remark}

% =====================================================================
\section{T4: reconfiguration is catalytic and not extractable}
\label{sec:catalyst}
% =====================================================================

We show that a reconfiguration is not an element of the graph it edits: it is
neither a vertex nor a path, but a change in the reachability relation, and
cannot be recovered as a located object from the reconfigured graph alone.

\begin{definition}[Catalytic edit]
\label{def:catalytic}
A reconfiguration $\reshuffle:\Gph\mapsto\Gph'$ is \emph{catalytic for the
cell $U$} if it changes the reachability $\pi(\cdot,U)$ but is not itself a
vertex of $\V$ nor a single path of $\Gph$ or $\Gph'$; formally,
$\reshuffle$ is the difference of reachability structures,
$\reshuffle=\pi'(\cdot,U)-\pi(\cdot,U)$, a function on vertices, not an
element of $\V$ or a walk.
\end{definition}

\begin{theorem}[Reconfigurations are not extractable]
\label{thm:notextract}
Let $\reshuffle$ be a catalytic edit carrying $\Gph$ to $\Gph'$. There is no
vertex $v\in\V'$ and no single path $P$ in $\Gph'$ such that $\reshuffle$
equals $v$ or $P$; the edit is recoverable only as the \emph{relation}
$\pi'-\pi$, and two distinct graphs $\Gph_1\neq\Gph_2$ may admit the same
$\reshuffle$ \textup{(}the same change in reachability\textup{)}. Hence
$\reshuffle$ cannot be located inside $\Gph'$ as one of its parts.
\end{theorem}

\begin{proof}
By \Cref{def:catalytic}, $\reshuffle$ is a function $\V\to\Z_{\ge0}$
(the per-vertex change in edge-disjoint route count), an object of a
different type from a vertex (an element of $\V$) or a path (a walk in
$\Gph'$); no element of one type equals an element of another. That distinct
graphs may share a $\reshuffle$: take $\Gph_1,\Gph_2$ differing by an
automorphism of the introduced route; both realise the same
$\pi'-\pi$, so $\reshuffle$ does not determine the graph and is not a
located part of it. Therefore $\reshuffle$ is recoverable only as the
reachability difference, not as a vertex or path of $\Gph'$.
\end{proof}

\begin{corollary}[Reached from many, present in none]
\label{cor:manynone}
If a cell $\Inv$ is an invariant (multiply-reachable), the reconfiguration
that made it multiply-reachable is reached by---i.e.\ consistent with---many
routes and is identical to none of them. It is the change in the route
structure, not a route.
\end{corollary}

\begin{proof}
Multiple reachability means $\ge2$ edge-disjoint routes
(\Cref{def:invariant}); the edit raising $\pi$ is, by \Cref{thm:notextract},
the reachability difference, not any one of those routes. Hence it is
consistent with all and equal to none.
\end{proof}

\begin{remark}[Reading]
\label{rem:catalyst-read}
Read: an idea is a \emph{catalyst}. It reshuffles which routes reach a truth
without being a route or a vertex, is not consumed into the result, and
cannot be pointed at inside the reconfigured graph---only its effect on the
routes can. This is why ``you cannot extract the idea'': the idea is the
reconfiguration, and a reconfiguration is a relation-change, not a located
part. No theorem depends on the reading.
\end{remark}

% =====================================================================
\section{T5: every finite specification carries slack}
\label{sec:slack}
% =====================================================================

A reconfiguration may be \emph{specified} by a finite description---a finite
list of edits. We show that any such description admits a non-identity
re-description, so the description is never identical to the reconfiguration
it specifies.

\begin{definition}[Specification; slack]
\label{def:spec}
A \emph{specification} of a reconfiguration $\reshuffle$ is a finite sequence
of edits whose net effect is $\reshuffle$. A specification has \emph{slack}
if there is a distinct specification with the same net effect $\reshuffle$
(a different list realising the same reachability change), or a removable
edit whose deletion leaves $\reshuffle$ unchanged.
\end{definition}

\begin{theorem}[Specifications carry slack]
\label{thm:slack}
Let $\reshuffle$ be a reconfiguration introducing at least one route, and let
$\Sigma$ be any finite specification of $\reshuffle$ in a graph with
$\abs{\V}\ge3$. Then $\Sigma$ has slack: there exists a distinct
specification $\Sigma'\neq\Sigma$ with the same net effect, or a removable
edit. Consequently no specification is uniquely determined by its effect;
the map (specification)$\to$(reconfiguration) is non-injective.
\end{theorem}

\begin{proof}
By \Cref{thm:notextract}, distinct graphs/edits may realise the same
$\reshuffle$; concretely, the introduced route may be re-routed through an
alternative edge-disjoint path (which exists once $\abs{\V}\ge3$ and the
target is multiply-reachable, \Cref{def:invariant}), or padded by an edit and
its inverse-in-effect, yielding $\Sigma'\neq\Sigma$ with the same net
reachability change. Hence $\Sigma$ is not the unique realiser of
$\reshuffle$; the specification carries slack and is not identical to
$\reshuffle$.
\end{proof}

\begin{corollary}[The minimal specification is not the reconfiguration]
\label{cor:notthething}
Even a minimal specification (one with no removable edit) is one realiser
among others of $\reshuffle$; it is a particular route to the
reconfiguration, not the reconfiguration, which is the reachability change
common to all its realisers.
\end{corollary}

\begin{proof}
By \Cref{thm:slack} the effect $\reshuffle$ has $\ge2$ realisers whenever a
re-route exists; a minimal specification is one of them. The invariant
content---$\reshuffle$ itself---is what they share, not what any one is.
\end{proof}

\begin{remark}[Reading]
\label{rem:slack-read}
Read: if it can be put into words (a finite specification), it is not the
idea but a route to the idea, because some other words would have served
(slack). The idea is the reachability change common to all the sentences
that specify it; each sentence is a particular realiser with removable
choices, and the existence of an alternative realiser is the proof that the
sentence is not the idea. No theorem depends on this reading.
\end{remark}

% =====================================================================
\section{T6: the four modes and the bridge}
\label{sec:modes}
% =====================================================================

We classify reconfigurations by two binary invariants and isolate a mode
whose endpoints close holonomy while its route does not.

\begin{definition}[Two invariants of a reconfiguration]
\label{def:twoinv}
For a reconfiguration $\reshuffle$ introducing a route to a target cell $U$,
define:
\begin{enumerate}[label=\textup{(\alph*)},leftmargin=2.2em]
\item \emph{Closure} $\kappa(\reshuffle)\in\{0,1\}$: $\kappa=1$ iff the
introduced route's target cell $U$ is an invariant whose cycles all have
$\hol=0$ (holonomy-closing); $\kappa=0$ otherwise.
\item \emph{Survival} $s(\reshuffle)\in\{0,1\}$: $s=1$ iff after
$\reshuffle$ the target retains $\pi(\cdot,U)\ge2$ (alternate routes
survive); $s=0$ iff $\reshuffle$ leaves or makes $U$ singly-reachable
(alternates removed).
\end{enumerate}
\end{definition}

\begin{theorem}[Four modes]
\label{thm:fourmodes}
Every route-introducing reconfiguration falls in exactly one of four classes
by $(\kappa,s)$:
\[
\begin{array}{ll}
(\kappa,s)=(1,1): & \text{adds a holonomy-closing route; alternates survive.}\\
(\kappa,s)=(0,1): & \text{adds a non-closing route; alternates survive.}\\
(\kappa,s)=(1,0): & \text{adds a closing route while removing alternates.}\\
(\kappa,s)=(0,0): & \text{adds a non-closing route while removing alternates.}
\end{array}
\]
The classes are disjoint and exhaustive, and each is realisable.
\end{theorem}

\begin{proof}
$(\kappa,s)\in\{0,1\}^2$ partitions reconfigurations into four disjoint,
exhaustive classes by definition. Realisability: for $(1,1)$ add an
edge-disjoint route to a multiply-reachable invariant; for $(0,1)$ add a
route to a non-invariant cell while leaving $\pi\ge2$ elsewhere; for $(1,0)$
add a closing route and delete competing routes until $\pi=1$; for $(0,0)$
do the latter to a non-closing target. Each construction satisfies
\Cref{ax:finite} on $\Gph'$.
\end{proof}

\begin{definition}[Bridge]
\label{def:bridge}
A \emph{bridge} is a reconfiguration whose introduced route has
holonomy-closing \emph{endpoints} (each stated terminal cell is an invariant
with $\hol=0$) but passes through an intermediate cell $U^\dagger$ that is
\emph{not} an invariant and whose local cycles have $\hol\neq0$. Formally:
the route $u\rightsquigarrow U^\dagger\rightsquigarrow U$ has
$\kappa$-closing $u$ and $U$ but a non-closing $U^\dagger$ on the path.
\end{definition}

\begin{theorem}[Endpoint-audit cannot detect a bridge]
\label{thm:bridge}
Let $\reshuffle$ be a bridge with stated endpoints $u,U$ and hidden
intermediate $U^\dagger$. Then:
\begin{enumerate}[label=\textup{(\roman*)},leftmargin=2.2em]
\item \textbf{Endpoints pass.} Every cycle confined to the stated endpoint
cells has $\hol=0$: an audit restricted to the endpoints detects no
inconsistency.
\item \textbf{Route fails.} There is a cycle through $U^\dagger$ with
$\hol\neq0$: an audit of the route detects the inconsistency.
\item \textbf{Undetectable by endpoints alone.} No function of the endpoint
holonomies determines $\hol$ at $U^\dagger$; detection requires evaluating
holonomy along the introduced route, not at its termini.
\end{enumerate}
\end{theorem}

\begin{proof}
(i) By hypothesis the endpoint cells are invariants with vanishing cycle
holonomy, so any cycle within them sums to $0$. (ii) By hypothesis some cycle
through $U^\dagger$ has $\hol\neq0$; it exists by \Cref{def:bridge}. (iii)
The endpoint holonomies are all $0$ regardless of $\hol(U^\dagger)$, so they
carry no information about it; $\hol(U^\dagger)$ is a separate cycle sum,
recoverable only by traversing the route through $U^\dagger$. Hence no
endpoint-only function determines it.
\end{proof}

\begin{corollary}[Route-audit detects the bridge]
\label{cor:routeaudit}
A bridge is detected if and only if holonomy is evaluated along its
introduced route (through $U^\dagger$), not merely at its stated endpoints.
The minimal detector is the cycle through $U^\dagger$.
\end{corollary}

\begin{proof}
Immediate from \Cref{thm:bridge}(ii)--(iii).
\end{proof}

\begin{remark}[Reading, on which nothing depends]
\label{rem:modes-read}
The four modes read, respectively, as: a reconfiguration that adds a true,
checkable route while leaving other checks in place; one that adds a
plausible route to an un-anchored cell while leaving checks in place; one
that adds a true route while removing the other checks; and one that adds an
un-anchored route while removing the checks. A \emph{bridge} reads as a
reconfiguration assembled from true, checkable endpoints whose
\emph{intermediate} cell is not checkable: every stated claim closes its
cycle, yet a cycle through the unstated intermediate does not.
\Cref{thm:bridge} reads ``checking only the stated endpoints cannot catch
it; only auditing the route can.'' We assert only the graph statements; the
readings are orientation and no theorem depends on them.
\end{remark}

\begin{remark}[A neutrality reading]
\label{rem:neutral}
Note that a single introduced route, considered in isolation, carries no
$(\kappa,s)$ value: closure and survival are properties of the route
\emph{within the whole graph} (which cell it reaches, whether cycles close,
whether alternates remain), not of the route as a sequence of edges. The
same edge-sequence may be class $(1,1)$ in one graph and part of a bridge in
another. A reader may take this as the statement that the moral or epistemic
status of a step is a property of the path-in-its-graph, not of the step;
identical local edits differ only in the routes they are steps of. No theorem
depends on this reading.
\end{remark}

% =====================================================================
\section{The Master Theorem}
\label{sec:master}
% =====================================================================

\begin{theorem}[Forced by finiteness with a floor]
\label{thm:master}
The consequences T0--T6 hold for every finite contact graph
\textup{(}\Cref{ax:finite}\textup{)}, and fail to be forced when finiteness
or the positive floor is dropped:
\begin{enumerate}[label=\textup{(\alph*)},leftmargin=2.2em]
\item \textbf{Zero floor.} If $\inf_e\wt(e)=0$, a cell may be separated at
arbitrarily small cost, so T0 fails; a route may be introduced below any
positive content, so the unit T3 is not floored; and an invariant may be
approached at vanishing cost, weakening T1.
\item \textbf{Infinite vertex set.} If $\abs{\V}=\infty$, a vertex may have
infinite degree, so a minimal route-introducing edit need not exist (T3),
and reachability counts $\pi$ need not be finite, weakening T1 and T4.
\end{enumerate}
\end{theorem}

\begin{proof}
The forward direction is the conjunction of
\Cref{thm:floor,thm:region,thm:holonomy,thm:unit-floor,thm:nonreturn,%
thm:notextract,thm:slack,thm:fourmodes,thm:bridge}, each proved under
\Cref{ax:finite}. (a): a graph with cuts of weight $2^{-n}\to0$ violates the
conclusion of \Cref{thm:floor}; the rest follow as stated. (b): an infinite
star has a centre of infinite degree, defeating the minimality and finiteness
used in \Cref{thm:unit-floor,thm:notextract}. Hence finiteness with a
positive floor is necessary as well as sufficient.
\end{proof}

% =====================================================================
\section{Discussion: the dictionary}
\label{sec:dict}
% =====================================================================

We collect, in one place and as orientation only, the correspondence between
the graph structures proved above and the everyday vocabulary they were
designed to replace. \emph{No theorem depends on any entry of this
dictionary.} Each everyday term names exactly the graph structure beside it
and carries no further commitment.

\begin{center}
\begin{tabular}{ll}
\toprule
\textbf{Graph structure} & \textbf{Reading (orientation only)}\\
\midrule
Vertex / proper cell & a distinguished part / a resolved thing\\
Edge weight $\ge\floorc$ & cost of telling two parts apart\\
Floor $\floorc>0$ (T0) & no distinction is free\\
Invariant cell (T1) & a truth: multiply-reached, region not point\\
Cycle holonomy $\hol$ (T2) & whether routes to a truth agree\\
$\hol=0$ on all cycles & consistent (``anchored'') reachability\\
Reconfiguration unit (T3) & an idea, in unit form\\
Monotone count, non-return (T3) & a received idea cannot be un-received\\
Catalytic edit, not extractable (T4) & an idea is a catalyst, not a thing\\
Specification slack (T5) & if it can be said, it is not the idea\\
$(\kappa,s)=(1,1)$ & adding a checkable route, checks intact\\
$(\kappa,s)=(0,\cdot)$ & adding an un-anchored route\\
$(\kappa,s)=(\cdot,0)$ & removing the alternate checks\\
Bridge (T6) & true endpoints, un-anchored hidden route\\
Route-audit (T6 cor.) & checking the path, not the stated endpoints\\
\bottomrule
\end{tabular}
\end{center}

\begin{remark}[Why the dictionary is severed]
\label{rem:severed}
The everyday terms are contested and their informal definitions invite
dispute. By proving every result of the abstract graph and confining the
vocabulary to this table, we ensure that disagreement about the words cannot
touch the theorems: one may reject every reading in the right-hand column and
the left-hand column stands unchanged. The mathematical content is the
graph theory; the readings are an interpretation a reader may take or leave.
\end{remark}

% =====================================================================
\section{Conclusion}
\label{sec:conclusion}
% =====================================================================

From a single premise---a finite graph whose edges cost something---we
derived a determinate chain about how its reachable structure changes under
minimal committed edits. A separating cut is never free (T0). A
multiply-reached, relabelling-stable cell is a region fixed by its
complement, never a point (T1). A group-valued transport has a cycle
holonomy that vanishes exactly on consistent reachability (T2). The minimal
committed edit introduces one floor-crossing route, increments a monotone
count, and is therefore irreversible (T3). Such an edit is a catalyst---a
change in the reachability relation, neither vertex nor path, not extractable
from the result (T4). Every finite specification of it carries removable
slack, so no specification is identical to what it specifies (T5). Edits
classify by closure and survival into four modes, among which a
\emph{bridge}---true endpoints, un-anchored intermediate---passes every
endpoint check and is caught only by auditing the route (T6). All are forced
by finiteness with a floor and fail without it (Master Theorem).

The everyday correspondents---thing, knowledge, truth, idea, propagation,
and the modes of true and false transmission---are recorded in the
dictionary of \Cref{sec:dict} as orientation on which no theorem depends.
The results are statements about finite weighted graphs, and stand or fall
as such.

% =====================================================================
\bibliographystyle{plainnat}
\begin{thebibliography}{9}
\bibitem[Diestel(2017)]{diestel2017} Reinhard Diestel.
\emph{Graph Theory}, 5th edition. Graduate Texts in Mathematics 173,
Springer, 2017.
\bibitem[Bollob\'as(1998)]{bollobas1998} B\'ela Bollob\'as.
\emph{Modern Graph Theory}. Graduate Texts in Mathematics 184,
Springer, 1998.
\end{thebibliography}

\end{document}
